\documentclass[12pt,a4paper]{article}

\usepackage{fontspec}
\setmainfont{Libertinus Serif}
\usepackage[greek.ancient,spanish]{babel}
\babelfont[greek]{rm}{Libertinus Serif}
\usepackage[a4paper,margin=3.2cm]{geometry}
\usepackage{microtype}
\usepackage[hidelinks]{hyperref}

% Griego politónico escrito directamente en UTF-8:
% \gr{...} para palabras sueltas; bloques con \foreignlanguage{greek}{...}
\newcommand{\gr}[1]{\foreignlanguage{greek}{#1}}
\newcommand{\stanzabreak}{\par\vspace{0.6\baselineskip}}

\setlength{\parindent}{1.25em}
\setlength{\parskip}{0.35em}
\raggedbottom

\usepackage{amssymb}
\usepackage{enumitem}

\hypersetup{
  pdftitle={Inventario editorial de citas: Leer la Odisea en tiempos iletrados},
  pdfauthor={}
}

\title{Inventario editorial de citas\\[0.4em]\large\emph{Leer la Odisea en tiempos iletrados}}
\author{}
\date{Agosto de 2026}

\begin{document}

\maketitle

\section*{Objeto}

Este documento identifica los textos ajenos reproducidos en el manuscrito activo y propone un orden práctico para revisar atribuciones y permisos. No determina por sí mismo qué usos están legalmente amparados: esa decisión corresponde a la editorial y, cuando proceda, a sus asesores.

El inventario excluye los pasajes comentados en los archivos LaTeX. Las cantidades de versos son aproximadas, pero suficientes para establecer prioridades.

\section*{Prioridad inmediata}

\subsection*{1. Traducción de Juan Manuel Rodríguez Tobal}

La traducción de la \emph{Odisea} publicada por Hiperión en 2026 es la fuente principal de las citas homéricas. El manuscrito reproduce aproximadamente 2.380 versos distribuidos por casi todos los capítulos, además de varios pasajes del prólogo en «Una osada Odisea».

No conviene gestionar estos fragmentos uno por uno. La acción útil es solicitar una autorización global que identifique con claridad:

\begin{itemize}[leftmargin=*]
  \item la obra de origen y su edición;
  \item la extensión aproximada reproducida;
  \item la edición impresa y digital del ensayo;
  \item los territorios e idiomas previstos;
  \item la fórmula de crédito que debe aparecer en el libro;
  \item si la autorización corresponde al traductor, a Hiperión o a ambos según el contrato de edición.
\end{itemize}

Este permiso condiciona el libro entero y debe resolverse antes que cualquier cita secundaria.

\subsection*{2. Kavafis en traducción de José María Álvarez}

Hay dos usos de la traducción publicada por Hiperión:

\begin{itemize}[leftmargin=*]
  \item «Deslealtad»: veintiséis versos en «Helena»;
  \item «Ítaca»: tres versos en «Marcharse de Ítaca».
\end{itemize}

La atribución ya figura en el texto, pero conviene confirmar edición, derechos de la traducción y fórmula de crédito. Puesto que también interviene Hiperión en la traducción de Tobal, puede intentarse una gestión editorial conjunta.

\subsection*{3. Lluís Llach}

«Marcharse de Ítaca» reproduce tres versos en catalán de \emph{Viatge a Ítaca}. Aunque procedan de Kavafis, el texto cantado por Llach constituye una versión distinta de la traducción castellana de José María Álvarez. Deben identificarse y consultarse por separado los derechos y el crédito de esa versión.

\section*{Citas que necesitan identificación o comprobación}

\subsection*{4. Tolkien}

La introducción reproduce cinco versos del poema del Anillo en castellano. El texto menciona la edición de Minotauro comprada en 1982, pero no identifica al traductor de esos versos ni la edición exacta. Antes de solicitar o descartar un permiso, hay que comprobar esos dos datos en el ejemplar utilizado.

\subsection*{5. Poesía española}

La introducción cita tres versos de Miguel Hernández y tres de Federico García Lorca. La autoría está indicada en la prosa, pero faltan el título del poema y la procedencia editorial.

\begin{itemize}[leftmargin=*]
  \item Miguel Hernández: «Tristes armas / si no son las palabras. / Tristes. Tristes», de \emph{Cancionero y romancero de ausencias}.
  \item Federico García Lorca: «Te acordaste de mí, cuando subías, / al silencio que sufre la serpiente / prisionera de grillos y de umbrías», de «El poeta pregunta a su amor por la ciudad encantada de Cuenca».
\end{itemize}

La editorial debe verificar la situación aplicable a cada texto y decidir si basta la atribución o si procede alguna gestión adicional.

\subsection*{6. Otras traducciones clásicas}

El manuscrito incluye fragmentos de varias traducciones castellanas, además de la de Tobal:

\begin{itemize}[leftmargin=*]
  \item tres versos de la \emph{Eneida} traducida por Javier de Echave-Sustaeta;
  \item quince versos de la \emph{Odisea} de José Manuel Pabón;
  \item dos fragmentos breves en prosa de la \emph{Odisea} de Luis Segalá y Estalella.
\end{itemize}

Las atribuciones aparecen en el manuscrito. Queda por fijar una referencia bibliográfica homogénea y comprobar con la editorial si la situación de cada traducción exige alguna actuación.

\section*{Textos cedidos o recibidos personalmente}

\subsection*{7. Candelas Gala}

«Circe» reproduce dos pasajes del epílogo de Candelas Gala a \emph{Abandonando Ítaca}. Aunque se trate de un libro del propio autor, el texto citado es de Candelas Gala. Conviene dejar constancia escrita de su autorización y confirmar si Eolas debe intervenir por la edición de origen.

\subsection*{8. Juan Manuel Rodríguez Tobal y Quique Ruiz}

«Una osada Odisea» reproduce una respuesta en verso de Juan Manuel Rodríguez Tobal y el poema completo «Multiverso», de Quique Ruiz. Ambos textos llegaron en el marco de una comunicación personal. La acción necesaria es sencilla: conservar una autorización escrita para reproducirlos y acordar la forma exacta del crédito.

El capítulo incluye además una frase de un correo de Tobal sobre el efecto sonoro de \emph{territodaseñor todaterritremendo}. Por prudencia editorial, puede incluirse en la misma autorización.

\section*{Material que no plantea una gestión de traducción ajena}

\begin{itemize}[leftmargin=*]
  \item Los poemas de \emph{Abandonando Ítaca} y «Lecciones de conducir» son del autor.
  \item Las traducciones que el manuscrito identifica expresamente como propias ---varios fragmentos de la \emph{Ilíada} y las glosas del griego--- no requieren permiso de otro traductor.
  \item El texto griego de Homero puede reproducirse como original clásico; debe conservarse, no obstante, una pauta uniforme para indicar canto y versos cuando resulte útil.
  \item Las menciones de títulos, nombres y frases muy breves no forman parte de este inventario, salvo cuando integran una cita poética reconocible.
\end{itemize}

\section*{Orden de trabajo recomendado}

\begin{enumerate}[leftmargin=*]
  \item $\square$ Resolver la autorización global de Tobal/Hiperión.
  \item $\square$ Tramitar con Hiperión los dos poemas de Kavafis traducidos por José María Álvarez.
  \item $\square$ Comprobar la versión de Llach y la traducción del poema de Tolkien.
  \item $\square$ Homogeneizar las referencias de Echave-Sustaeta, Pabón y Segalá.
  \item $\square$ Documentar las autorizaciones personales de Candelas Gala, Quique Ruiz y Juan Manuel Rodríguez Tobal.
  \item $\square$ Decidir con la editorial la forma final de los créditos: notas al pie, página de créditos o ambas.
\end{enumerate}

\section*{Conclusión}

La cuestión decisiva es una sola: la autorización para reproducir extensamente la traducción de Tobal. Resuelta esa gestión, el resto del inventario se reduce a un grupo pequeño de citas identificables y a varias autorizaciones personales fáciles de documentar.

No se propone modificar todavía el manuscrito. La forma de los créditos debe decidirse después de conocer las condiciones de cada autorización, para evitar añadir notas que luego haya que rehacer.

\end{document}
